\documentclass[11pt]{article}

\usepackage[letterpaper,margin=1in]{geometry}
\usepackage[T1]{fontenc}
\usepackage{lmodern}
\usepackage{amsmath}
\usepackage{amssymb}
\usepackage{enumitem}
\usepackage{booktabs}
\usepackage[hidelinks]{hyperref}

\setlist[enumerate]{itemsep=1pt,topsep=3pt}

\newcommand{\cN}{\mathcal N}
\newcommand{\cK}{\mathcal K}
\newcommand{\cG}{\mathcal G}
\newcommand{\cP}{\mathcal P}
\newcommand{\cA}{\mathcal A}
\newcommand{\TR}{\mathit{TR}}

\title{Solving the Simplified Model, Step by Step}
\author{}
\date{}

\begin{document}
\maketitle

\noindent This document derives the solution of the simplified model, one substitution
at a time. It adds no assumptions and no results to the model section. Its purpose is
to show how the equilibrium conditions reduce to a small system that can be solved,
which objects have closed forms, and which ones must be computed. Propositions and
theorems cited by name and number, and all notation, are those of the model section,
whose notation table lists every symbol.

The model has quantity competition among groups, Cobb--Douglas preferences across
sectors, labor as the only input, non-negative tariffs, and plants that serve a given
set of markets. Section~\ref{sec:entry} adds the extension in which groups choose which
plants serve each market.

\tableofcontents

\setcounter{section}{-1}
\section{What Has to Be Solved}
\label{sec:overview}

\paragraph{Primitives.} The economy is described by the sets of countries $\cN$ and
sectors $\cK$; labor endowments $L_n$; spending shares $\beta_k$ and elasticities
$\sigma_k>1$; the groups, with their headquarters, sectors and plants; plant
productivities $\varphi_i$; the markets each plant serves; trade costs $d_{\ell n}$;
tariffs $t_{\ell nk}\ge0$; and ownership shares $\theta_{gn}$. All of these are given.

\paragraph{Unknowns.} An equilibrium consists of wages $w_n$, incomes $X_n$, and the
price $p_{i,n}$ and quantity $q_{i,n}$ of every plant in every market it serves.
Market sizes $E_{nk}$, price indices $P_{nk}$, shares, markups and profits follow from
these.

\paragraph{Strategy.} There is no closed form for the equilibrium as a whole. Every
object, however, can be written as a function of the wage vector $w$ alone, through a
chain of substitutions in which each link is either a closed form or a problem with a
unique and easily computed solution. The equilibrium is then a wage vector that
satisfies $N-1$ labor-market conditions.

\begin{center}
\small
\begin{tabular}{@{}llll@{}}
\toprule
Step & Solves for & Takes as given & How \\
\midrule
1 & demand functions & prices, spending & closed form \\
2 & unit and delivered costs & wages & closed form \\
3 & shares, markups, prices in one market & delivered costs & one scalar equation \\
4 & market sizes and incomes & wages & linear system \\
5 & labor-market conditions & everything above & closed form \\
6 & wages & --- & $N-1$ equations \\
\bottomrule
\end{tabular}
\end{center}

\noindent Steps 1--3 hold at any wages. Market shares do not depend on spending, so
Step 3 can be solved before Step 4, and incomes then solve a linear system.

\section{Step 1: Household Demand}
\label{sec:demand}

\paragraph{Within a sector.} Fix country $n$ and sector $k$, and write $\sigma=\sigma_k$.
To obtain an amount $C_{nk}$ of the sector's bundle at least cost, the household
chooses $q_{i,n}$ to minimize $\sum_{i\in\cA_{nk}}p_{i,n}q_{i,n}$ subject to
$\bigl[\sum_iq_{i,n}^{(\sigma-1)/\sigma}\bigr]^{\sigma/(\sigma-1)}=C_{nk}$. With
multiplier $\chi$, the first-order condition for variety $i$ is
\[
  p_{i,n}=\chi\,\frac{\partial C_{nk}}{\partial q_{i,n}}
  =\chi\,C_{nk}^{1/\sigma}q_{i,n}^{-1/\sigma},
  \qquad\text{so}\qquad
  q_{i,n}=\Bigl(\frac{p_{i,n}}{\chi}\Bigr)^{-\sigma}C_{nk}.
\]
Substituting into the constraint,
$C_{nk}=\chi^{\sigma}C_{nk}\bigl[\sum_ip_{i,n}^{1-\sigma}\bigr]^{\sigma/(\sigma-1)}$,
so $\chi=\bigl[\sum_ip_{i,n}^{1-\sigma}\bigr]^{1/(1-\sigma)}\equiv P_{nk}$. The
multiplier is the price index: the cost of one unit of the bundle. Spending on the
sector is then $P_{nk}C_{nk}$, and
\begin{equation}
  q_{i,n}=p_{i,n}^{-\sigma}P_{nk}^{\sigma-1}\,(P_{nk}C_{nk}).
  \label{eq:lowertier}
\end{equation}

\paragraph{Across sectors.} The household maximizes $U_n=\prod_kC_{nk}^{\beta_k}$
subject to $\sum_kP_{nk}C_{nk}=X_n$. With multiplier $\zeta$, the first-order condition
for sector $k$ is $\beta_kU_n/C_{nk}=\zeta P_{nk}$, so $P_{nk}C_{nk}=\beta_kU_n/\zeta$.
Adding over sectors and using $\sum_k\beta_k=1$ gives $U_n/\zeta=X_n$, hence
\begin{equation}
  E_{nk}\equiv P_{nk}C_{nk}=\beta_kX_n .
  \label{eq:spending}
\end{equation}
Each sector receives a fixed share of income. Substituting into \eqref{eq:lowertier},
\begin{equation}
  q_{i,n}=p_{i,n}^{-\sigma}P_{nk}^{\sigma-1}E_{nk},
  \qquad
  s_{i,n}\equiv\frac{p_{i,n}q_{i,n}}{E_{nk}}=\Bigl(\frac{p_{i,n}}{P_{nk}}\Bigr)^{1-\sigma}.
  \label{eq:demand}
\end{equation}
The revenue shares $s_{i,n}$ add up to one in every market. Substituting the optimal
bundles into utility gives $U_n=X_n/P_n$ with $P_n=\prod_k(P_{nk}/\beta_k)^{\beta_k}$,
which is why welfare is real income.

\section{Step 2: Costs}
\label{sec:costs}

\paragraph{Unit cost.} A plant located in $\ell$ produces $y=\varphi_il$ units with $l$
units of labor. Its cost of producing one unit is
\begin{equation}
  a_i=\frac{w_\ell}{\varphi_i}.
  \label{eq:unitcost}
\end{equation}

\paragraph{Delivered cost.} To sell $q_{i,n}$ units in $n$, the plant produces
$d_{\ell n}q_{i,n}$ units, because of iceberg trade costs. Its cost of serving $n$,
exclusive of tariffs, is
\begin{equation}
  V_{i,n}=a_id_{\ell n}q_{i,n},
  \label{eq:V}
\end{equation}
country $n$ collects $t_{\ell nk}V_{i,n}$ in tariffs, and the cost of delivering one
more unit to $n$, tariff included, is
\begin{equation}
  c_{i,n}=a_i\,d_{\ell n}\,(1+t_{\ell nk})
  =\frac{w_\ell\,d_{\ell n}(1+t_{\ell nk})}{\varphi_i}.
  \label{eq:delivered}
\end{equation}

\paragraph{Where the money goes.} The whole of $V_{i,n}$ is paid to labor in $\ell$.
Out of the plant's revenue in $n$, what is left after $V_{i,n}$ and the tariff is profit
(Step 3).

\section{Step 3: One Market}
\label{sec:market}

Fix a market $(n,k)$ and drop its indices: $\sigma$, $E$, $P$, $c_i$, $s_i$ and
$t_i\equiv t_{\ell(i)nk}$. Let $T_g$ be the plants of group $g$ in the market. The inputs
of this step are the delivered costs $c_i$, which by \eqref{eq:delivered} depend only on
wages, and the spending level $E$. The outputs are prices, quantities, shares, markups
and profits.

\subsection{Inverse demand}

Let $\gamma=(\sigma-1)/\sigma$ and $Q=(\sum_iq_i^\gamma)^{1/\gamma}$, the quantity of
the sector's bundle. Since $PQ=E$, \eqref{eq:lowertier} can be inverted as
$p_i=P(q_i/Q)^{-1/\sigma}=(E/Q)(q_i/Q)^{-1/\sigma}$, so that
\begin{equation}
  p_i=E\,Q^{-\gamma}q_i^{\gamma-1},
  \qquad
  s_i=\frac{p_iq_i}{E}=\frac{q_i^\gamma}{Q^\gamma} .
  \label{eq:inverse}
\end{equation}

\subsection{The group's first-order condition and the markup}

Group $g$ chooses $q_i$ for every $i\in T_g$ to maximize
$\pi_g=\sum_{i\in T_g}(p_i-c_i)q_i$. Its revenue is
\[
  \sum_{j\in T_g}p_jq_j=E\,Q^{-\gamma}\sum_{j\in T_g}q_j^\gamma=E\,S_g,
  \qquad
  S_g\equiv\sum_{j\in T_g}s_j .
\]
Using $\partial\ln Q/\partial q_i=q_i^{\gamma-1}/Q^\gamma=s_i/q_i$, the marginal revenue of
plant $i$ is
\[
  \frac{\partial(ES_g)}{\partial q_i}
  =E\,Q^{-\gamma}\Bigl[\gamma\,q_i^{\gamma-1}-\gamma\,\frac{q_i^{\gamma-1}}{Q^\gamma}\sum_{j\in T_g}q_j^\gamma\Bigr]
  =\gamma\,p_i\,(1-S_g).
\]
The first term is the revenue from selling more. The second is cannibalization:
selling more lowers the price index and takes sales from every variety in the market,
including the group's other plants, and the group counts the loss on all of them, which
is proportional to its total share $S_g$. Setting marginal revenue equal to marginal
cost $c_i$ gives
\begin{equation}
  p_i=\mu(S_g)\,c_i,
  \qquad
  \mu(S)=\frac{\sigma}{(\sigma-1)(1-S)} .
  \label{eq:markup}
\end{equation}
The markup is the same for every plant of the group and depends only on the group's
total share. It starts at $\mu(0)=\sigma/(\sigma-1)$, the markup of a negligible firm,
and is unbounded as $S\to1$, which is why every market needs at least two groups.

\subsection{Shares within a group}

With a common markup, $p_i/p_j=c_i/c_j$ for any two plants of the same group, so by
\eqref{eq:demand} $s_i/s_j=(c_i/c_j)^{1-\sigma}$. Adding over the group,
\begin{equation}
  \frac{s_i}{S_g}=\frac{c_i^{1-\sigma}}{K_g},
  \qquad
  K_g\equiv\sum_{i\in T_g}c_i^{1-\sigma}.
  \label{eq:withingroup}
\end{equation}
The group gives its cheaper plants a larger part of its sales. Costs show up in how
the group splits its sales, not in its markups.

\subsection{The share equation}

Let $A\equiv\sum_ip_i^{1-\sigma}=P^{1-\sigma}$, so that $s_i=p_i^{1-\sigma}/A$.
Substituting $p_i=\mu(S_g)c_i$ and adding over the group,
\[
  S_g=\sum_{i\in T_g}\frac{\mu(S_g)^{1-\sigma}c_i^{1-\sigma}}{A}
  =\frac{\mu(S_g)^{1-\sigma}K_g}{A},
\]
that is,
\begin{equation}
  \psi(S_g)=\frac{K_g}{A},
  \qquad
  \psi(S)\equiv S\,\mu(S)^{\sigma-1}
  =\Bigl(\frac{\sigma}{\sigma-1}\Bigr)^{\sigma-1}\frac{S}{(1-S)^{\sigma-1}} .
  \label{eq:sharemap}
\end{equation}
For a given $A$ this is one equation in one unknown for each group. Taking logs,
\begin{equation}
  \ln S_g-(\sigma-1)\ln(1-S_g)=\ln K_g-\ln A-(\sigma-1)\ln\frac{\sigma}{\sigma-1}
  \equiv x_g(A).
  \label{eq:shareeq}
\end{equation}
The left side rises from $-\infty$ to $+\infty$ as $S_g$ goes from $0$ to $1$, so the
equation has exactly one root $S_g(A)\in(0,1)$ for every $A>0$. It has no closed form in
general, but Newton's method converges quickly from any interior start, using the
derivative of the left side, $[1+(\sigma-2)S]/[S(1-S)]$.

\subsection{The market-clearing equation}

Shares must add up to one, so the market's equilibrium is the value of $A$ that
solves
\begin{equation}
  \mathcal S(A)\equiv\sum_gS_g(A)=1 .
  \label{eq:clearing}
\end{equation}
Differentiating \eqref{eq:sharemap} in logs, with
$d\ln\psi/d\ln S=1+(\sigma-1)S/(1-S)$,
\begin{equation}
  \frac{dS_g}{d\ln A}=-S_g\Lambda_g,
  \qquad
  \Lambda_g\equiv\frac{1-S_g}{1+(\sigma-2)S_g}>0 .
  \label{eq:dSdA}
\end{equation}
Every share falls as $A$ rises, so $\mathcal S$ is strictly decreasing and
\eqref{eq:clearing} has at most one solution (Proposition~2). A bracket that contains it
is found as follows. Since $\mu(S)\ge\mu(0)$, \eqref{eq:sharemap} gives
$S_g\le K_g\mu(0)^{1-\sigma}/A$, so at $A_{\mathrm{hi}}=\mu(0)^{1-\sigma}\sum_gK_g$ shares
add up to at most one. As $A\to0$ every share tends to one, so $\mathcal S$ tends to the
number of groups, which is at least two; halving $A$ from $A_{\mathrm{hi}}$ reaches a
point $A_{\mathrm{lo}}$ with $\mathcal S(A_{\mathrm{lo}})\ge1$ in finitely many steps.
Between the two, Newton's method on $\ln A$, with derivative
$\mathcal S'(\ln A)=-\sum_gS_g\Lambda_g$ and a bisection step whenever a Newton step
would leave the bracket, converges to the unique root.

\subsection{Recovering everything else}

Once $A$ and the shares $S_g$ are known, every other object in the market follows in
closed form:
\begin{align}
  \mu_g&=\mu(S_g),\qquad p_i=\mu_gc_i,\qquad P=A^{1/(1-\sigma)},
  \label{eq:recover1}\\
  s_i&=S_g\,\frac{c_i^{1-\sigma}}{K_g},\qquad r_i=s_iE,\qquad q_i=\frac{r_i}{p_i},
  \label{eq:recover2}\\
  V_i&=\frac{r_i}{\mu_g(1+t_i)},\qquad \text{tariff revenue } t_iV_i .
  \label{eq:recover3}
\end{align}
The price index is consistent with the prices, since
$\sum_ip_i^{1-\sigma}=\sum_g\mu_g^{1-\sigma}K_g=A\sum_gS_g=A$. Profit is
\[
  \pi_g=\sum_{i\in T_g}(p_i-c_i)q_i=\Bigl(1-\frac{1}{\mu_g}\Bigr)\sum_{i\in T_g}r_i
  =E\,S_g\,m(S_g),
  \qquad
  m(S)\equiv1-\frac{1}{\mu(S)}=\frac{1+(\sigma-1)S}{\sigma},
\]
so that, with $\kappa=\sigma-1$,
\begin{equation}
  \pi_g=\frac{E}{\sigma}\bigl(S_g+\kappa S_g^2\bigr).
  \label{eq:profit}
\end{equation}

\paragraph{The revenue split.} For each plant, revenue splits exactly into labor cost,
tariff and profit:
\begin{equation}
  r_i=\mu_g(1+t_i)V_i=\underbrace{V_i}_{\text{labor}}
  +\underbrace{t_iV_i}_{\text{tariff}}
  +\underbrace{(\mu_g-1)(1+t_i)V_i}_{\text{profit}} .
  \label{eq:revid}
\end{equation}
Dividing by $E$ and adding over plants, each dollar spent in the market splits as
\begin{equation}
  \underbrace{\sum_i\frac{s_i}{\mu_{g(i)}(1+t_i)}}_{v:\ \text{labor}}
  +\underbrace{\sum_i\frac{t_is_i}{\mu_{g(i)}(1+t_i)}}_{\tau:\ \text{tariffs}}
  +\underbrace{\sum_gS_g\,m(S_g)}_{\text{profit}}=1 .
  \label{eq:perdollar}
\end{equation}

\paragraph{Two properties used below.} First, shares, markups and the price index
depend only on delivered costs, not on $E$: spending scales quantities, revenue, costs
and profits, and nothing else. Second, multiplying every delivered cost by the same
factor leaves every share and markup unchanged and multiplies every price and $P$ by
that factor. Since delivered costs are proportional to wages, scaling all wages scales
all prices and leaves all shares and markups unchanged.

\section{Step 4: Market Sizes and Incomes}
\label{sec:incomes}

From now on wages $w$ are given. By Steps 2 and 3, every delivered cost, share and
markup is then known, and everything that remains is linear in spending.

\paragraph{Profits and tariffs as a function of spending.} Group $g$ of sector $k$
earns $E_{n'k}S_{g,n'}m(S_{g,n'})$ in each destination $n'$ it serves, and country $n$
collects $\TR_n=\sum_k\tau_{nk}E_{nk}$ in tariffs, where $\tau_{nk}$ is the tariff part
of \eqref{eq:perdollar} in market $(n,k)$. With $\sum_n\theta_{gn}=1$, the income of
country $n$ is
\[
  X_n=w_nL_n+\sum_g\theta_{gn}\sum_{n'}E_{n'k(g)}S_{g,n'}m(S_{g,n'})+\sum_k\tau_{nk}E_{nk}.
\]

\paragraph{Substituting market sizes.} By \eqref{eq:spending}, $E_{n'k}=\beta_kX_{n'}$.
Substituting,
\begin{equation}
  X_n=w_nL_n+\sum_{n'\in\cN}\Gamma_{nn'}X_{n'},
  \qquad
  \Gamma_{nn'}\equiv\sum_k\beta_k\Bigl[\sum_{g:\,k(g)=k}\theta_{gn}S_{g,n'}m(S_{g,n'})
  +\mathbf 1\{n'=n\}\,\tau_{nk}\Bigr].
  \label{eq:incomesys}
\end{equation}
$\Gamma_{nn'}$ is the income that country $n$ receives, as profit or tariff revenue, out
of each dollar that country $n'$ spends. In matrix form,
\begin{equation}
  X=wL+\Gamma X
  \qquad\Longrightarrow\qquad
  X=(I-\Gamma)^{-1}wL .
  \label{eq:XofW}
\end{equation}

\paragraph{Why the solution exists, is unique and exceeds labor income.} Every entry of
$\Gamma$ is non-negative, because profits and tariffs are. The column of $\Gamma$ for
country $n'$ adds up, over the receiving countries $n$, to
\[
  \sum_n\Gamma_{nn'}=\sum_k\beta_k\Bigl[\sum_gS_{g,n'}m(S_{g,n'})+\tau_{n'k}\Bigr]
  =1-\sum_k\beta_kv_{n'k}<1,
\]
by \eqref{eq:perdollar}: of each dollar country $n'$ spends, the part that goes to labor
never comes back as profit or tariff revenue. A non-negative matrix whose columns add
up to less than one has all its eigenvalues below one in absolute value, so
\[
  (I-\Gamma)^{-1}=I+\Gamma+\Gamma^2+\cdots\ge0,
  \qquad
  X=wL+\Gamma wL+\Gamma^2wL+\cdots\ge wL .
\]
The series has a simple reading. Labor income $wL$ is spent; part of that spending
returns as profits and tariffs, $\Gamma wL$; that income is spent again, and so on. Every
country's income is at least its labor income (Proposition~3). Market sizes follow from
\eqref{eq:spending}, and with them every quantity, cost and profit through
\eqref{eq:recover2}--\eqref{eq:profit}.

\paragraph{Scale.} Scaling all wages by a common factor leaves $\Gamma$ unchanged and
scales $wL$, so it scales incomes and market sizes by the same factor.

\section{Step 5: The Labor Market}
\label{sec:labor}

\paragraph{The wage bill.} Country $n$'s wage bill is the labor cost of all plants
located in $n$, in every market they serve. Let
\begin{equation}
  v_{n,(n',k)}\equiv\sum_{i\in\cA_{n'k}:\,\ell(i)=n}\frac{s_{i,n'}}{\mu_{g(i),n'}(1+t_{nn'k})}
  \label{eq:vcoef}
\end{equation}
be the share of spending in market $(n',k)$ that is paid to labor in $n$. Then, using
\eqref{eq:spending},
\begin{equation}
  W_n(w)=\sum_{n'\in\cN}\sum_{k\in\cK}v_{n,(n',k)}\,\beta_k\,X_{n'} .
  \label{eq:wagebill}
\end{equation}
Every term is known from Steps 3 and 4. The excess demand for labor is
\begin{equation}
  z_n(w)=\frac{W_n(w)}{w_n}-L_n,
  \label{eq:excess}
\end{equation}
and labor markets clear when $z_n=0$ for every $n$.

\paragraph{Walras' law.} One of these conditions is redundant, because at any wages
\[
  \sum_nw_nz_n(w)=\sum_nW_n-\sum_nw_nL_n=0 .
\]
To see it, add the revenue split \eqref{eq:revid} over all plants and markets: world
spending equals the world wage bill plus world tariff revenue plus world profits,
$\sum_nX_n=\sum_nW_n+\sum_n\TR_n+\sum_g\pi_g$. Adding the income equations over
countries, with $\sum_n\theta_{gn}=1$, gives
$\sum_nX_n=\sum_nw_nL_n+\sum_g\pi_g+\sum_n\TR_n$. The two expressions are equal, so
$\sum_nW_n=\sum_nw_nL_n$.

\paragraph{The numeraire.} Scaling all wages by a common factor scales $W_n$ and $w_n$
by that factor (Steps 3 and 4), so $z$ does not change. Only relative wages are
determined, and we set $w_1=1$. By Walras' law, $z_1=0$ whenever $z_2=\dots=z_N=0$.

\section{Step 6: Wages}
\label{sec:wages}

\paragraph{The reduced system.} Steps 2--5 reduce the whole model to $N-1$ equations in
$N-1$ unknowns:
\begin{equation}
  z_n(1,w_2,\dots,w_N)=0,\qquad n=2,\dots,N .
  \label{eq:reduced}
\end{equation}
Written out, the condition for country $n$ is
\[
  w_nL_n=\sum_{n'}\sum_k\beta_kX_{n'}
  \sum_{i\in\cA_{n'k}:\,\ell(i)=n}\frac{s_{i,n'}}{\mu_{g(i),n'}(1+t_{nn'k})},
  \qquad
  X=(I-\Gamma)^{-1}wL,
\]
where shares and markups depend on wages only through the delivered costs
$c_{i,n'}=w_{\ell(i)}d_{\ell(i)n'}(1+t_{\ell(i)n'k})/\varphi_i$. Evaluating $z(w)$ at a
given $w$ therefore means: compute delivered costs, solve every market (Step 3), solve
the income system (Step 4), and compute the wage bills (Step 5).

\paragraph{Three channels.} A rise in a foreign wage $w_j$ affects country $n$'s wage
bill in three ways. Plants located in $j$ become more expensive, so plants located in
$n$ gain sales, both within their own groups and against rival groups. The share
changes move markups, which changes how much of each dollar of revenue is paid to labor.
And incomes $X_{n'}$ change, because they include the profits of the groups each
country owns. The last channel is the ownership channel of the model. Theorem~1 states
when the first two always favor country $n$ and the third does not undo them.

\paragraph{Damped tâtonnement.} The algorithm updates wages by
\begin{equation}
  w_n^{\text{new}}=w_n^{1-\lambda}\Bigl(\frac{W_n(w)}{L_n}\Bigr)^{\lambda},
  \qquad\text{that is,}\qquad
  \ln w_n^{\text{new}}=\ln w_n+\lambda\ln\Bigl(1+\frac{z_n(w)}{L_n}\Bigr),
  \label{eq:tatonnement}
\end{equation}
with $\lambda\in(0,1)$, and then rescales so that $w_1=1$, which is allowed because the
update is homogeneous of degree one. A country's wage rises when its labor is in excess
demand, and a wage vector is a fixed point exactly when $W_n=w_nL_n$ for every $n$.

\paragraph{How much damping.} Let $M_{nj}=\partial\ln W_n/\partial\ln w_j$. In logs, the
update has derivative $(1-\lambda)I+\lambda M$. Its off-diagonal entries are positive
when labor is a gross substitute, and its diagonal entries are positive when
$\lambda<1/(1-M_{nn})$ for every $n$. When labor demand is elastic, $M_{nn}<0$, the bound
is below one, and an undamped update ($\lambda=1$) can overshoot.

\paragraph{Newton steps.} Close to the solution, Newton's method on
$F_n(\ln w)=\ln W_n-\ln w_n-\ln L_n$, $n=2,\dots,N$, converges faster. Its derivative is
$\partial F_n/\partial\ln w_j=M_{nj}-\mathbf 1\{n=j\}$, for $n,j\ge2$. The matrix $M$ is
obtained column by column, by changing one $\ln w_j$ slightly and re-evaluating Steps
2--5. Iterations stop when $\max_n|z_n|/L_n$ falls below a tolerance.

\section{The Complete Algorithm, and What to Check}
\label{sec:summary}

From the innermost step outward:
\begin{enumerate}[label=(\Alph*)]
\item \emph{Costs.} Given $w$, compute $a_i$ and $c_{i,n}$ by
\eqref{eq:unitcost}--\eqref{eq:delivered}.
\item \emph{Markets.} In every market, solve each share equation \eqref{eq:shareeq} for
$S_g(A)$ and the clearing equation \eqref{eq:clearing} for $A$; recover $\mu_g$, $p_i$,
$P$ and $s_i$ by \eqref{eq:recover1}--\eqref{eq:recover2}.
\item \emph{Incomes.} Build $\Gamma$ and compute $X$ by \eqref{eq:XofW} and $E$ by
\eqref{eq:spending}.
\item \emph{Labor market.} Compute $W_n$ by \eqref{eq:wagebill} and $z_n$ by
\eqref{eq:excess}.
\item \emph{Wages.} Update $w$ by \eqref{eq:tatonnement}, then by Newton steps,
repeating (A)--(D), until $z(w)=0$.
\end{enumerate}

\begin{center}
\small
\begin{tabular}{@{}ll@{}}
\toprule
Closed form & Computed \\
\midrule
demand, unit and delivered costs & each group's share, given $A$ (one monotone equation) \\
markups, prices, quantities given shares & each market's $A$ (one monotone equation) \\
shares within a group, profits & incomes (one linear system) \\
revenue split into labor, tariffs, profit & wages ($N-1$ equations) \\
wage bills given incomes and shares & \\
\bottomrule
\end{tabular}
\end{center}

\paragraph{Checks at the solution.} Each of these holds at an equilibrium and is a
useful test of any implementation.
\begin{enumerate}
\item In every market, shares add up to one and $\sum_ip_i^{1-\sigma}=P^{1-\sigma}$.
\item For every plant, revenue equals labor cost plus tariff plus profit,
\eqref{eq:revid}.
\item Every country's income equals labor income plus owned profits plus tariff
revenue, and is at least its labor income.
\item Every labor market clears, and Walras' law holds also away from the solution,
which tests the accounting at arbitrary wages.
\item The hypotheses of Theorem~1: no group holds more than half of any market, and
$\bar E^{(j)}_n>0$ for every pair of countries.
\end{enumerate}

\section{Extension: Endogenous Entry}
\label{sec:entry}

With entry, each group chooses which of its potential plants $\cP_g$ serve each market,
and a plant that serves market $(n,k)$ pays the fixed cost $F_{i,n}=f_kw_\ell
f^d_{\ell n}$ in labor of its location. Steps 1--3 are unchanged for given plant sets.
What changes is that plant sets are now unknowns.

\subsection{Entry into one market}

\paragraph{The value of a plant set.} By Step 3, a set $T\subseteq\cP_g$ affects the
group's profit only through its strength $K(T)=\sum_{i\in T}c_i^{1-\sigma}$, and it
costs $\Phi(T)=\sum_{i\in T}F_{i,n}$. Taking $A$ as given, the group's share with set $T$
is $\psi^{-1}(K(T)/A)$, and its net payoff is
\begin{equation}
  U_g(T;A)=\Omega\bigl(K(T);A\bigr)-\Phi(T),
  \qquad
  \Omega(K;A)=\frac{E}{\sigma}\bigl(S+\kappa S^2\bigr)\Big|_{S=\psi^{-1}(K/A)} .
  \label{eq:payoff}
\end{equation}
Serving the market with no plant gives $U_g=0$.

\paragraph{The ladder.} A group with $|\cP_g|$ potential plants has $2^{|\cP_g|}$ sets.
Since $\Omega$ rises with $K$, a set can be chosen only if no other set has at least as
much strength at no greater fixed cost. The construction is:
\begin{enumerate}
\item compute $(K(T),\Phi(T))$ for every set $T$;
\item discard every $T$ for which some $T'$ has $K(T')\ge K(T)$ and
$\Phi(T')\le\Phi(T)$, with one inequality strict;
\item order the remaining sets by strength; fixed cost rises along the order too. Call
them rungs $0,1,\dots,R_g$, with rung 0 the empty set.
\end{enumerate}
When all of a group's plants pay the same fixed cost $F$ in the market, rung $r$ is the
set of the $r$ plants with the lowest delivered costs. A stand-alone firm has two rungs:
out or in.

\paragraph{The best response.} At a given $A$, the group picks the rung with the highest
payoff,
\begin{equation}
  K^*_g(A)=K\bigl(T^*_g(A)\bigr),
  \qquad
  T^*_g(A)=\arg\max_{r=0,\dots,R_g}\;\Omega(K_r;A)-\Phi_r .
  \label{eq:bestresponse}
\end{equation}
With equal fixed costs this is a cutoff: the strength added by each rung falls as more
expensive plants are added, and $\Omega$ is strictly concave in $K$ (Theorem~2(a)), so
the gain from one more rung falls, and the group climbs while the gain is positive.

\paragraph{The entry equilibrium of the market.} Substituting the best responses into
\eqref{eq:clearing} gives one equation in $A$:
\begin{equation}
  \mathcal S^*(A)\equiv\sum_g\psi^{-1}\!\Bigl(\frac{K^*_g(A)}{A}\Bigr)=1 .
  \label{eq:entryclearing}
\end{equation}
When no group would hold more than half of the market with any of its rungs, each
$K^*_g(A)$ does not rise with $A$ (Theorem~2), so $\mathcal S^*$ falls strictly with $A$:
strength does not rise and the aggregate is higher. It crosses one at most once. It is solved by
bisection on $\ln A$.

\paragraph{The integer problem.} Because $K^*_g(A)$ moves in jumps, $\mathcal S^*$ can
jump over one. The bisection then converges to the aggregate $\bar A$ where some group
switches rungs. If the group is assigned its larger rung, the clearing aggregate lies at
or above $\bar A$, where the smaller rung is its best response; if it is assigned its
smaller rung, the clearing aggregate lies below $\bar A$, where the larger rung is.
Either way the group wishes to deviate, and no exact equilibrium in whole plants exists.
The procedure keeps the configuration with the fewest groups wishing to deviate and
reports how many there are.

\subsection{General equilibrium with entry}

\paragraph{What changes in Steps 4 and 5.} Fixed costs are paid to labor where the plant
is located and are deducted from the group's profit. Let $\Phi_g$ be the total fixed
cost of group $g$'s active plants, which depends on wages and plant sets only. Income
becomes
\[
  X=wL-\theta'\Phi+\Gamma X
  \qquad\Longrightarrow\qquad
  X=(I-\Gamma)^{-1}\bigl(wL-\theta'\Phi\bigr),
\]
where $\theta'\Phi$ is the vector with entries $\sum_g\theta_{gn}\Phi_g$, and each
country's wage bill \eqref{eq:wagebill} adds the fixed costs of the plants located in
it. Walras' law still holds, because every fixed cost is both a deduction from profit
and a payment to labor. Scaling all wages scales fixed costs too, so it leaves every
entry decision unchanged.

\paragraph{The outer loop.} Plant sets and wages are solved by an outer loop:
\begin{enumerate}
\item start from a configuration of plant sets, for example every potential plant
active in every market of its sector;
\item given the plant sets, solve Steps 2--6 for wages, and record in every market the
delivered costs, fixed costs, spending and aggregate $A$;
\item in every market, compute each group's ladder and best response at the recorded
aggregate, and its gain from switching;
\item if no group gains, stop: every market is in an entry equilibrium at the
equilibrium wages;
\item otherwise, let the groups with the largest gains switch, and return to 2.
\end{enumerate}
Moving only the groups with the largest gains damps the loop. Moving every group at once
can make the configuration oscillate, because wages computed for many active plants
differ from wages for few. The choice affects the path of the algorithm, not its fixed
points, and imposes no order of entry. The loop is not guaranteed to converge; when it
does not reach a configuration in which no group wishes to move, it returns the one with
the fewest such groups. Because the joint equilibrium of wages and plant sets is not
known to be unique, the procedure is run from several starting points.

\end{document}
